\documentclass[paper=a4, twoside, fontsize=12pt, parskip=half-, numbers=enddot, open=right]{scrbook}
\usepackage{preambule}
\usepackage{macros}
\begin{document}

\chapter*{Feuille de synthèse \\ Justifications des angles (chapitre 5)}

\begin{center}
   \emph{Pour rappel, à chaque angle sa propre justification ! } 
\end{center}


\subsubsection{Structurer les justifications se fait de cette façon : 
}

\begin{center}
    \begin{tblr}{
            column{1}={3cm},
            column{2}={7cm},
            rows={t},
        }
        $|\widehat{Angle}|= \dots \degres$ & \textbf{CAR} justification justification justification justification justification justification justification justification justification justification justification justification justification justification justification justification justification justification justification justification justification justification justification  \\
        $|\widehat{Angle}|= \dots \degres$ & \textbf{CAR}  justification justification justification justification justification justification justification justification justification justification justification justification justification justification justification justification justification justification justification justification justification justification justification
    \end{tblr}
\end{center}

\subsubsection{Quelques justifications possibles }
\begin{enumerate}[itemsep=.7cm, label=-]
    \item \emph{[insérer]} et \emph{[insérer]} sont des angles supplémentaires \\ + calcul 
    \item \emph{[insérer]} et \emph{[insérer]} sont des angles complémentaires \\ + calcul 
    \item \emph{[insérer]} et \emph{[insérer]} sont des angles opposés par le sommet
    \item Les angles \emph{[insérer]} et \emph{[insérer]} sont des angles correspondants et \emph{[insérer]} $\sslash$ \emph{[insérer]}
    \item Les angles \emph{[insérer]} et \emph{[insérer]} sont des angles alternes internes et \emph{[insérer]} $\sslash$   \emph{[insérer]}
    \item Les angles \emph{[insérer]} et \emph{[insérer]} sont des angles alternes externes et \emph{[insérer]} $\sslash$  \emph{[insérer]}
    \item Les angles d'un triangle équilatéral valent toujours $60 \degres$
    \item La somme des amplitudes des angles d'un triangle vaut $180 \degres$ \\ + calcul
    \item La somme des amplitudes des angles d'un quadrilatère vaut $360 \degres$ \\ + calcul
    \item La somme des amplitudes des angles d'un \emph{[insérer un polygone à $n$ côtés]} vaut \emph{[calcul à l'aide de la formule]} \\ + calcul
    \item La \emph{[insérer]} est la bissectrice de l'angle \emph{[insérer]}, elle coupe donc cet angle en deux parts égales \\ + calcul 
    \item Les angles à la base d'un triangle isocèle ont même amplitude 
    \item La somme des amplitudes des angles d'un triangle vaut $180 \degres$ et les angles à la base d'un triangle isocèle ont même amplitude \\ + calcul 
    \item \ldots
\end{enumerate}

\end{document}
